\documentclass[12pt]{article}
\usepackage[utf8]{inputenc}
\usepackage[T1]{fontenc}
\usepackage{hyperref}
\title{Zenaga Language Resources for K-12 Curriculum Development}
\author{Compiled by Claude Sonnet 5.5 (claude-sonnet-5-5) for Andr\'as Kornai}
\date{}
\begin{document}
\maketitle

\section{Introduction}
Zenaga is a highly divergent Berber (Afro-Asiatic) language of southern Mauritania. Sources found in search give widely differing speaker estimates, and UNESCO is reported to rate it severely endangered. No school materials or native-language teaching programmes turned up in search. The documentation found is scholarly and mostly in French.

\section{Primary Educational Resources}
No primers, textbooks or school programmes in Zenaga were found in search.

\section{Linguistic and Reference Materials}
\begin{itemize}
\item Catherine Taine-Cheikh, \emph{Dictionnaire z\'enaga-fran\c{c}ais} (Tome 1) and \emph{Dictionnaire fran\c{c}ais-z\'enaga} (Tome 2), published by K\"oppe. The introduction contains a phonological and morphosyntactic study. \url{https://www.koeppe.de/titel_dictionnaire-zenaga-francais-tome-1-dictionnaire-francais-zenaga-tome-2}
\item Bookseller listing for the French-Zenaga dictionary: \url{https://abebooks.com/9783896459275/Dictionnaire-fran%C3%A7ais-z%C3%A9naga-berb%C3%A8re-Mauritanie-3896459279/plp}
\item Overview article: \url{https://en.wikipedia.org/wiki/Zenaga_language}
\item Numeral data for Zenaga (Trarza region), Max Planck Institute for Evolutionary Anthropology: \url{https://lingweb.eva.mpg.de/channumerals/Zenaga.htm}
\end{itemize}

\section{Religious and Community Resources}
No religious or community-organisation resources for Zenaga were found in search.

\section{Digital and Audio Resources}
No audio, video or software resources were found in search.

\section{Curriculum Development Context}
Zenaga has very few speakers, and the evidence points to older adults as the main first-language speakers. Any curriculum would have to start from documentation, building on the Taine-Cheikh dictionary and grammar, and would need work with the remaining speakers. Little exists beyond scholarly description.
\end{document}